\section{Criterios de evaluación y líneas base}
\label{sec:met_protocolo}

\subsection{Selección de modelos, métricas y criterio de decisión}
\label{sec:met_seleccion}

Las confusiones entre categorías no tienen el mismo costo clínico. Confundir MEL con NV retrasa la
extirpación de una lesión cuya letalidad depende del estadio en que se diagnostica; confundir NV con MEL
lleva a extirpar una lesión benigna sin necesidad; las confusiones entre las clases benignas tienen menor
consecuencia~\cite{tschandl2018ham10000,nci2026melanoma}. Esa asimetría obliga a medir el desempeño clase
por clase, además de en conjunto, y a dar a las clases minoritarias el mismo peso que a la mayoritaria.

\pendiente{El Capítulo~\ref{cp:marco} define qué mide cada métrica; esta subsección indica cuáles se usan
para decidir y por qué. Se justifican la sensibilidad macro como métrica principal de clasificación y el F1
macro como complemento, el umbral de especificidad $\geq 0.85$ fijado en validación para la detección de
melanoma, las métricas de calibración y las dos de generación (Nivel 1 con referencia; Nivel 2 con
\emph{clinical entity recall} y \emph{diagnostic consistency}), con la distinción entre fidelidad diagnóstica
y validez clínica.

Se describen a continuación los barridos de hiperparámetros y el criterio de selección, que opera
enteramente en validación: el test no interviene en ninguna decisión. Hay que distinguir dos niveles. El
\emph{checkpoint} de cada corrida es la época con mejor F1 ponderado en validación. La corrida ganadora
entre las del barrido, en cambio, se elige por \emph{F1 macro} en validación, porque el F1 ponderado está
dominado por la clase mayoritaria ---NV es el 67\,\% de las imágenes--- y la sensibilidad macro, que sería
la alternativa natural por ser la métrica principal de este trabajo, ignora los falsos positivos y por lo
tanto premia a los modelos que sobre-predicen las clases minoritarias. El F1 macro es el único de los tres
que pondera las clases por igual y a la vez penaliza los falsos positivos. La sensibilidad macro se reporta
como resultado principal, pero no se usa para elegir.

Los resultados se reportan como media $\pm$ desvío entre semillas.}
